\chapter{Background}
\label{ch:background}

This chapter collects the material that later chapters rely on. \Cref{sec:bg:quant} reviews
integer quantization and the effect of a single bit flip on an INT8 weight.
\Cref{sec:bg:npu} outlines the software stack through which a quantized model reaches a
Qualcomm NPU. \Cref{sec:bg:transformer} describes the weight layout of a decoder-only
transformer. \Cref{sec:bg:crypto} reviews the cryptographic primitives used for attestation,
and \cref{sec:bg:stats} the statistics used by the activation monitor.

% ---------------------------------------------------------------------------
\section{Integer quantization}
\label{sec:bg:quant}

\subsection{Affine and symmetric quantization}

Quantization maps a real-valued tensor $x$ to integers $q$ using a scale $s>0$ and a zero
point $z_0$ \cite{jacob2018quantization,nagel2021whitepaper}:
\begin{equation}
  q = \clamp\!\Big(\round\big(x/s\big) + z_0,\; q_{\min},\; q_{\max}\Big),
  \qquad \hat x = s\,(q - z_0).
  \label{eq:affine}
\end{equation}
For signed 8-bit storage, $q_{\min}=-128$ and $q_{\max}=127$. \emph{Symmetric} quantization
fixes $z_0=0$ and chooses $s = \max|x| / 127$, so that the largest-magnitude value maps to
$\pm127$. \emph{Asymmetric} quantization uses unsigned storage $[0,255]$ with
$s=(x_{\max}-x_{\min})/255$ and $z_0 = \round(-x_{\min}/s)$, which uses the full range
when the distribution is skewed. \textsc{EdgeGuard}'s helper \code{compute\_int8\_params}
implements both variants per tensor (\cref{ch:implementation}).

The reconstruction error of a value inside the clipping range is at most $s/2$. Rounding errors
are commonly modelled as uniform on $[-s/2,s/2]$ with variance $s^2/12$. This matters for the
monitor in two ways. Activations observed at runtime are integers, so their empirical variance
includes this quantization noise. And statistics computed in FP32 at calibration time and
compared against INT8 values at runtime would be measured in different units unless they are
rescaled. Calibrating directly on the INT8 values avoids this mismatch by construction.

\subsection{Per-tensor, per-channel and low-bit formats}

A single scale per tensor is simple but wastes resolution when channel magnitudes differ.
Per-channel quantization assigns a scale to each output channel of a weight matrix and is the
usual choice for convolution and linear weights \cite{nagel2021whitepaper}. Below eight bits,
formats such as INT4 pack two values per byte. \textsc{EdgeGuard} stores INT4 values in the
low and high nibbles of a byte and sign-extends them on unpacking. The attestation layer is
indifferent to these choices because it hashes the stored bytes, whatever they encode.

For LLMs, activation outliers make naive 8-bit activation quantization lossy. SmoothQuant
migrates part of the activation range into the weights with a mathematically equivalent
per-channel rescaling, enabling 8-bit weights and activations for models up to 530B parameters
\cite{xiao2023smoothquant}. GPTQ quantizes weights to 3--4 bits using approximate second-order
information \cite{frantar2023gptq}. Both produce integer tensors of the kind that
\textsc{EdgeGuard} attests.

\subsection{Two's complement and the cost of one bit}
\label{sec:bg:twos}

An INT8 value is stored as a byte $b_7b_6\dots b_0$ in two's complement:
\begin{equation}
  q = -128\,b_7 + \sum_{k=0}^{6} 2^k b_k .
\end{equation}
Flipping bit $k<7$ changes $q$ by exactly $\pm 2^k$. Flipping the sign bit $b_7$ changes $q$ by
$\pm128$ when interpreted in the signed range: for instance $-5$ (\code{0xFB}) becomes $123$
(\code{0x7B}), a change of $128$. In real units the change is $s\cdot 2^k$, so the
maximum single-flip perturbation of a weight is $128\,s$. \Cref{tab:bitflip_effect}
summarises the effect by bit position.

\begin{table}[t]
\centering
\caption{Effect of flipping one bit of an INT8 weight stored in two's complement.}
\label{tab:bitflip_effect}
\begin{tabular}{cccl}
\toprule
\textbf{Bit $k$} & \textbf{$|\Delta q|$} & \textbf{$|\Delta \hat x|$} & \textbf{Example ($q=-5$)} \\
\midrule
0 (LSB) & 1   & $s$     & $-5 \to -6$ \\
1       & 2   & $2s$    & $-5 \to -7$ \\
2       & 4   & $4s$    & $-5 \to -1$ \\
3       & 8   & $8s$    & $-5 \to -13$ \\
4       & 16  & $16s$   & $-5 \to -21$ \\
5       & 32  & $32s$   & $-5 \to -37$ \\
6       & 64  & $64s$   & $-5 \to -69$ \\
7 (sign)& 128 & $128s$  & $-5 \to 123$ \\
\bottomrule
\end{tabular}
\end{table}

A subtle implementation point follows from this arithmetic. In a high-level language, XORing a
\emph{signed} integer with $2^7$ does not reproduce the hardware effect. In Python,
\code{-5 \^{} 128} evaluates to $-133$, which is outside the INT8 range and either raises an
overflow error or, if clipped, saturates to $-128$ instead of wrapping to $123$. A faithful
simulator must operate on the unsigned byte view. Two such errors were found in the original
\textsc{EdgeGuard} simulators and are described in \cref{sec:impl:bugs}.

% ---------------------------------------------------------------------------
\section{The edge NPU software stack}
\label{sec:bg:npu}

\begin{figure}[t]
\centering
\begin{tikzpicture}[
  box/.style={draw=eggray, rounded corners=2pt, minimum width=10.5cm, minimum height=0.8cm,
              align=center, font=\small},
  hl/.style={box, fill=eglight, draw=egblue},
  node distance=0.18cm]
\node[box] (app) {Application (Android / Linux process)};
\node[box, below=of app] (fw) {Framework delegate: LiteRT, ONNX Runtime, ExecuTorch};
\node[hl, below=of fw] (qnn) {Qualcomm AI Engine Direct (QNN) API: backend, context, graph};
\node[box, below=of qnn, minimum width=3.3cm, xshift=-3.6cm] (cpu) {CPU backend};
\node[box, right=0.15cm of cpu, minimum width=3.3cm] (gpu) {GPU backend};
\node[box, right=0.15cm of gpu, minimum width=3.6cm, fill=eglight, draw=egblue] (htp) {HTP backend (NPU)};
\node[box, below=0.18cm of gpu, minimum width=10.5cm] (dram) {Shared DRAM: weights, activations, context binary};
\draw[-{Stealth}, eggray] (app) -- (fw);
\draw[-{Stealth}, eggray] (fw) -- (qnn);
\draw[-{Stealth}, eggray] (qnn.south) ++(-3.6,0) -- (cpu.north);
\draw[-{Stealth}, eggray] (qnn.south) -- (gpu.north);
\draw[-{Stealth}, eggray] (qnn.south) ++(3.45,0) -- (htp.north);
\node[right=0.3cm of dram, font=\scriptsize, text=egorange!80!black, align=left] {Rowhammer,\\cold boot};
\end{tikzpicture}
\caption{Simplified software stack for a quantized model on a Snapdragon device. The weights
that the NPU executes reside in DRAM shared with the rest of the system.}
\label{fig:npu_stack}
\end{figure}

Qualcomm's AI Engine Direct SDK, also referred to as QNN, provides a low-level API over the
Kryo CPU, Adreno GPU and Hexagon NPU, organised around tools, devices, backends, contexts,
graphs and an operation-package registry \cite{qualcomm_qnn}. Each processing unit is driven by a
\emph{backend} library; the HTP backend targets the Hexagon processor with its fused AI
accelerator \cite{qualcomm_qnn_backend}. Higher-level frameworks delegate to these backends.
ExecuTorch, for example, lowers PyTorch models to a QNN context binary for the HTP backend
\cite{executorch_qnn}.

Two properties of this stack matter here. First, the weights an NPU executes are stored in
system DRAM, which is shared with other software and exposed to DRAM-level attacks
(\cref{fig:npu_stack}). Second, the runtime is optimised for throughput and exposes inputs and
outputs, not intermediate activations, by default. Obtaining per-layer activations for
monitoring therefore requires either marking intermediate tensors as graph outputs at
compilation time or running an instrumented copy of the graph. \textsc{EdgeGuard}'s QNN adapter
is designed around the first option (\cref{sec:design:adapter}).

% ---------------------------------------------------------------------------
\section{Decoder-only transformer weight layout}
\label{sec:bg:transformer}

\begin{figure}[t]
\centering
\begin{tikzpicture}[
  w/.style={draw=egblue, fill=eglight, rounded corners=1.5pt, minimum height=0.55cm,
            font=\scriptsize\ttfamily, inner sep=2pt},
  n/.style={draw=eggray, rounded corners=1.5pt, minimum height=0.55cm, font=\scriptsize\ttfamily, inner sep=2pt},
  lab/.style={font=\scriptsize, text=eggray}]
\node[w, minimum width=4.2cm] (emb) at (0,0) {embed\_tokens  $V\times H$};
\begin{scope}[yshift=-1.0cm]
\node[n, minimum width=4.2cm] (ln1) at (0,0) {input\_layernorm  $H$};
\node[w, minimum width=2.0cm, anchor=north west] (q) at ($(ln1.south west)+(0,-0.15)$) {q\_proj $H\!\times\! H$};
\node[w, minimum width=2.0cm, anchor=north east] (k) at ($(ln1.south east)+(0,-0.15)$) {k\_proj $H\!\times\! d_h$};
\node[w, minimum width=2.0cm, anchor=north] (v) at ($(q.south)+(0,-0.15)$) {v\_proj $H\!\times\! d_h$};
\node[w, minimum width=2.0cm, anchor=north] (o) at ($(k.south)+(0,-0.15)$) {o\_proj $H\!\times\! H$};
\node[n, minimum width=4.2cm, anchor=north] (ln2) at ($(ln1.south)+(0,-1.6)$) {post\_attention\_layernorm $H$};
\node[w, minimum width=2.0cm, anchor=north west] (g) at ($(ln2.south west)+(0,-0.15)$) {gate\_proj $I\!\times\! H$};
\node[w, minimum width=2.0cm, anchor=north east] (u) at ($(ln2.south east)+(0,-0.15)$) {up\_proj $I\!\times\! H$};
\node[w, minimum width=4.2cm, anchor=north] (d) at ($(ln2.south)+(0,-0.85)$) {down\_proj $H\times I$};
\node[draw=eggray, dashed, rounded corners=3pt, inner sep=5pt, fit=(ln1)(q)(k)(v)(o)(ln2)(g)(u)(d)] (blk) {};
\node[lab, anchor=west] at (blk.east) {$\times\,N$ blocks (9 tensors each)};
\end{scope}
\node[n, minimum width=4.2cm, below=0.45cm of blk] (norm) {model.norm  $H$};
\node[w, minimum width=4.2cm, below=0.15cm of norm] (head) {lm\_head  $V\times H$};
\end{tikzpicture}
\caption{Weight tensors of the decoder-only transformer layout used by
\code{QuantizedTransformerStub}, with HuggingFace-style names. Shaded tensors are
matrices; outlined tensors are normalisation vectors. Grouped-query attention is represented by
the smaller key and value projections.}
\label{fig:transformer_layout}
\end{figure}

Llama~3, Qwen2.5 and Phi-3 are decoder-only transformers \cite{vaswani2017attention,
llama3herd,qwen25report,phi3report}. Each block applies RMS normalisation, multi-head
self-attention with query, key, value and output projections, a second normalisation and a
SwiGLU feed-forward network with gate, up and down projections. A token embedding matrix
precedes the blocks; a final normalisation and an output projection (``LM head'') follow them.

\textsc{EdgeGuard}'s \code{QuantizedTransformerStub} generates INT8 tensors with this layout
(\cref{fig:transformer_layout}). For hidden size $H$, feed-forward size $I$, key/value
projection width $d_h$, vocabulary $V$ and $N$ blocks, the stub has $9N+3$ tensors and
\begin{equation}
  P_{\text{stub}} = 2VH + N\big(2H^2 + 2Hd_h + 3HI + 2H\big) + H
  \label{eq:stub_params}
\end{equation}
parameters. The stub is a structural model, not a replica. It always uses a separate LM head
(some of the real models tie it to the embedding), sets the key/value width to a single head
dimension, and stores normalisation weights as INT8. Its parameter counts therefore differ from
the published model sizes, and this report always labels them as \emph{stub} counts.
\Cref{tab:profiles} lists the five profiles.

\begin{table}[t]
\centering
\caption{Stub profiles. Dimensions follow the named models; parameter counts are those of the
stub per \cref{eq:stub_params} and are \emph{not} the official model sizes.}
\label{tab:profiles}
\small
% AUTO-GENERATED
\begin{tabular}{lrrrrrr}
\toprule
\textbf{Profile key} & \textbf{Blocks} & \textbf{Hidden} & \textbf{FFN} & \textbf{Vocab} & \textbf{Tensors} & \textbf{Stub params} \\
\midrule
\texttt{llama-3.2-1b} & 16 & 2{,}048 & 8{,}192 & 128{,}256 & 147 & 1{,}469{,}122{,}560 \\
\texttt{llama-3.2-3b} & 28 & 3{,}072 & 8{,}192 & 128{,}256 & 255 & 3{,}452{,}611{,}584 \\
\texttt{phi-3.5-mini} & 32 & 3{,}072 & 8{,}192 & 32{,}064 & 291 & 3{,}235{,}974{,}144 \\
\texttt{qwen2.5-0.5b} & 24 & 896 & 4{,}864 & 151{,}936 & 219 & 627{,}387{,}264 \\
\texttt{tiny} & 4 & 64 & 256 & 1{,}000 & 39 & 366{,}144 \\
\bottomrule
\end{tabular}

\end{table}

% ---------------------------------------------------------------------------
\section{Cryptographic primitives}
\label{sec:bg:crypto}

\subsection{SHA-256}
SHA-256 is specified in FIPS~180-4 \cite{fips180-4}. It maps a message of arbitrary length to a
256-bit digest. The properties relied on here are collision resistance (it is infeasible to
find $x\neq x'$ with $\Hash(x)=\Hash(x')$) and second-preimage resistance (given $x$, it is
infeasible to find $x'\neq x$ with the same digest). No practical attack on either property is
known for SHA-256. Hashing is fast on commodity CPUs. On the evaluation machine, the
implementation in \cref{ch:implementation} achieved roughly \MerkThrMin{}--\MerkThrMax{}\,MiB/s over
weight tensors.

\subsection{Merkle trees}
A Merkle tree \cite{merkle1987digital} commits to an ordered list of values with a single
digest. Leaves are digests of the values, and each internal node is the digest of the
concatenation of its two children. The root commits to every leaf: changing any leaf changes
the root unless a collision is found. A membership proof for one leaf consists of the
$\lceil\log_2 n\rceil$ sibling digests on its path to the root. Merkle trees are used in
certificate transparency, version-control systems and blockchains, and they have a known
pitfall that is relevant to this project. If an odd node is paired with a copy of itself, two
different leaf lists can produce the same root, for example $[a,b,c]$ and $[a,b,c,c]$. The
pitfall is harmless here only because the verifier fixes the number of leaves; this point is
taken up in \cref{sec:design:merkle}.

\subsection{ECDSA over P-256}
The root is signed with the Elliptic Curve Digital Signature Algorithm on the NIST P-256 curve,
as specified in FIPS~186-5 \cite{fips186-5}. A signature over the 32-byte root lets a verifier
holding only the public key confirm that the root was produced by the holder of the private
key. In a deployment, the public key and the verification code would sit in a trusted
component, such as a TEE or a secure boot chain, so that an attacker who can rewrite DRAM
cannot also rewrite the reference.

% ---------------------------------------------------------------------------
\section{Statistics for the activation monitor}
\label{sec:bg:stats}

\subsection{Streaming mean and covariance}
Calibration observes $T$ activation vectors $\mathbf{a}_1,\dots,\mathbf{a}_T\in\R^d$ per
monitored layer. Storing them all is wasteful, and the na\"\i ve one-pass formula
$\frac1T\sum a^2 - (\frac1T\sum a)^2$ loses precision through cancellation. Welford's method
\cite{welford1962note} updates the mean and the sum of centred outer products incrementally:
\begin{align}
  \boldsymbol\delta_t &= \mathbf{a}_t - \boldsymbol\mu_{t-1}, &
  \boldsymbol\mu_t &= \boldsymbol\mu_{t-1} + \boldsymbol\delta_t / t, &
  M_t &= M_{t-1} + \boldsymbol\delta_t\,(\mathbf{a}_t - \boldsymbol\mu_t)^{\top},
  \label{eq:welford}
\end{align}
with $\hat\Sigma = M_T/(T-1)$. Each update costs $O(d^2)$ time and the state occupies $O(d^2)$
memory, independent of $T$.

\subsection{$z$-scores and the Mahalanobis distance}
For a new vector $\mathbf{a}$, the per-channel standard score is
$z_j = (a_j - \mu_j)/(\sigma_j+\epsilon)$. The Mahalanobis distance \cite{mahalanobis1936}
\begin{equation}
  D(\mathbf{a}) = \sqrt{(\mathbf{a}-\boldsymbol\mu)^\top \hat\Sigma^{+} (\mathbf{a}-\boldsymbol\mu)}
  \label{eq:mahal_bg}
\end{equation}
accounts for correlation between channels. Under a Gaussian model, $D^2$ follows a
$\chi^2_d$ distribution. Lee et al.\ used class-conditional Mahalanobis distances on hidden
features to detect out-of-distribution and adversarial inputs \cite{lee2018simple}. Using the
pseudo-inverse $\hat\Sigma^{+}$ keeps the computation defined when the covariance is singular,
which is common for INT8 activations because of repeated values.

\subsection{Multiple comparisons}
\label{sec:bg:multiple}
The monitor flags an input if \emph{any} of $m$ monitored channels exceeds a threshold. If each
channel exceeds a two-sided threshold $\tau$ with probability $p(\tau)=2(1-\Phi(\tau))$ on clean
data, and channels are independent, the probability that a clean input is flagged is
\begin{equation}
  \mathrm{FPR}(\tau) = 1 - \big(1 - p(\tau)\big)^{m}.
  \label{eq:fwer}
\end{equation}
This is the family-wise error rate familiar from multiple hypothesis testing. For $\tau=4$,
$p\approx\PSingleZFour$, which looks negligible for one channel, but with $m=\MChannels$
channels \cref{eq:fwer} gives \PredFPRZFour{}. The Bonferroni correction \cite{dunn1961multiple}
chooses a per-channel level of $\alpha/m$ to bound the family-wise rate by $\alpha$, giving the
threshold
\begin{equation}
  z^*(\alpha) = \Phi^{-1}\!\Big(1 - \frac{\alpha}{2m}\Big).
  \label{eq:bonf}
\end{equation}
Bonferroni is conservative when channels are positively correlated, but it requires no
estimate of the correlation structure. \Cref{sec:eval:fpr} shows that \cref{eq:fwer}
predicts the measured false-positive rate of the monitor closely.

\subsection{ROC analysis}
For a scalar detection statistic, the receiver operating characteristic plots the true-positive
rate against the false-positive rate as the threshold varies. The area under the curve (AUC)
equals the probability that a randomly chosen positive scores higher than a randomly chosen
negative. An AUC of 0.5 corresponds to chance. The evaluation reports ROC curves for the
maximum $|z|$ statistic.
